Gold und Silber haben je die Masse $m$ und die Dichten $\rho_\text{Au} = \rAuO$ und $\rho_\text{Ag} = \rAgO$ (FoTaBe S.~187). Wenn sich die Volumina beim Legieren einfach addieren, ist
\begin{align}
 \bar\rho &= \frac{2m}{V_\text{Au}+V_\text{Ag}} = \frac{2m}{\frac{m}{\rho_\text{Au}}+\frac{m}{\rho_\text{Ag}}} = \rRiF \\
   &= \frac{2\cdot\rAuO\cdot\rAgO}{\rAuO+\rAgO} \\
   &= \rRig \approx \result{\rRigP{0}{3}}
\end{align}
Der gemessene Wert kann abweichen, weil sich die Volumina beim Legieren nicht genau addieren: Die Atome der beiden Metalle ordnen sich in einem gemeinsamen Kristallgitter an und können dichter (oder weniger dicht) gepackt sein als in den reinen Metallen. Zudem können Hohlräume oder Lufteinschlüsse im Ring die gemessene Dichte verkleinern.
